\chapter{Comparability of the FBA's Unmatched Orders and the CDA Book}
\label{app:comparability}

A natural first instinct is to treat the FBA's set of not-yet-matched
orders at time $t$ as a stand-in for ``the CDA order book at $t$,'' and
then compute the usual book-based metrics (quoted spread, depth, book
imbalance) on it the same way they are computed on the CDA. This appendix
explains why that instinct does not survive scrutiny, and what the
metrics defined in \Cref{app:formulas} do instead.

\section{The Question}

A CDA order book at $t$ is a \emph{settled} object: every order that could
trade against the book already has, continuously, one order at a time.
Nothing sits in it by accident of timing. The FBA has no such object
between clears --- only a buffer of orders accumulating toward the next
batch, and, once that batch clears, whatever is left over. Two candidate
readings of ``the FBA's book at $t$'' exist, and they fail for different
reasons.

\section{Pre-Clearing: Settlement Mismatch and Inflated Depth}

Read immediately \emph{before} the batch containing $t$ clears, the
pending set mixes two things that must not be mixed. Using the notation
of \Cref{sec:fbanotation}: some pending orders are inframarginal at the
price the batch is about to clear at (every buy with $\pi_k > \clr$,
every sell with $\pi_k < \clr$) and will execute the instant the gate
closes --- they are not resting liquidity in any economic sense, only
flow that batching has temporarily delayed. Others are genuinely
marginal or out of the money and will remain unmatched. Counting the
whole pending set as if it were book depth overstates FBA liquidity
purely as an artifact of the batching delay, not because more liquidity
is actually standing there.

The deeper problem is one of kind, not convention. The CDA book at $t$
reflects a mechanism that has already finished matching everything it
can as of $t$. The FBA's pending set at the same $t$ is, by construction,
\emph{unsettled}: matching for the interval containing $t$ has not
happened yet. Comparing the two is comparing a post-matching state to a
pre-matching one, not two readings of the same kind of object taken at
different conventions.

\section{Post-Clearing: Lopsidedness and a Different Kind of ``Still There''}

Read immediately \emph{after} the batch containing $t$ clears, the
residual --- what rolls into the next batch --- fixes the settlement
problem: both objects are now ``what is left once everything that could
match, did.'' Two problems survive even so.

First, the residual is lopsided at the price that matters. Orders priced
away from $\clr$ on both sides roll over, so the leftover set usually has
a best buy and a best sell. At $\clr$ itself, however, it is typically
one-sided: the unexecuted residual share reported in \Cref{sec:res23} has
a median of exactly $1.00$ across the month, because most batches do not
cross at all, and even the batches that trade leave most of their longer
side unmatched. An imbalance read off this set describes which side the
last batch could not absorb, not the standing two-sided pressure that a
CDA book's imbalance describes.

Second, even where the residual is two-sided, the
\emph{reason} an order is still there differs between the two
mechanisms. A CDA resting order has been passed over continuously by
every unit of opposing flow that has arrived since it was posted --- its
survival is a standing, cumulative judgment by the market. An FBA
residual order is unmatched because of how much opposing volume happened
to land in that one specific batch window; it rolls forward and may
clear in the very next one. ``Still unmatched'' means something more
persistent in a CDA book than it does in a one-second FBA buffer.

\section{What the Metric Formulas Actually Use}

None of the above is a problem this thesis's metrics inherited silently;
each book-like FBA quantity in \Cref{app:formulas} was already defined to
sidestep one of the two issues above, for a reason now made explicit.

\begin{description}
  \item[Depth at best] uses $D(\clr)$/$S(\clr)$ --- the \emph{total}
  volume eligible to trade at the clearing price, matched and unmatched
  together, not the unmatched residual. This is the fair analogue of a
  CDA book's touch size (``how much was there at the best price''), and
  deliberately not an answer to ``how much is left over.'' It does not
  repeat the Pre-Clearing section's mixing problem, because it is
  evaluated \emph{after} $\clr$ is known: $D(\clr)$ and $S(\clr)$ are
  read off the one price that actually cleared, not off some provisional
  price guessed mid-accumulation, so there is no delayed-execution volume
  being miscounted as standing liquidity here.
  \item[Quoted spread] uses the post-clearing residual, and only for
  batches where both a best unfilled buy \emph{and} a best unfilled sell
  exist. Because unmatched orders on both sides roll over, this holds in
  almost every second, so the gate rarely binds; what remains is the
  second problem above, that the leftover set is not the same kind of
  object as a CDA book, which is why quoted spread is read with caution in
  \Cref{sec:metricreliability}.
  \item[Book imbalance] is not computed for the FBA at all. Unlike
  quoted spread, it has no gated, partial version here: a book-imbalance
  figure needs a resting two-sided book to read continuously, which, as
  argued above, the FBA simply does not have between clears, so this
  thesis does not build a proxy for an object that does not exist.
  \item[Almost everything else] --- executed volume, trade
  count, VWAP, effective and realized spread, price impact --- is built
  from the \emph{executed trade set}, not the unmatched set at all. A
  trade is timing-unambiguous for both mechanisms, since it only exists
  post-clearing either way (a CDA trade the instant of crossing, an FBA
  trade the instant the batch clears), which is exactly why trade-based
  metrics sidestep the pre/post-clearing dilemma above entirely rather
  than resolving it.
\end{description}

Put together, the FBA's unmatched orders --- the leftover residual this
appendix has been discussing --- appear in exactly one metric, quoted
spread, and only there under the two-sided gate above; depth at best
uses a related but distinct, already-cleared total instead, and book
imbalance is not attempted. Nowhere in this thesis is the unmatched set
read, on its own, as a general stand-in for ``the FBA order book.''
